\chapter{Breast Implants (Breast Augmentation)}

\section*{Definition and Overview}
Breast augmentation is a surgical procedure in which breast implants are placed to increase the size, shape, or fullness of the breasts. It is most commonly performed for cosmetic reasons, but it is also used to correct significant asymmetry, to restore volume lost through weight loss, pregnancy, or breastfeeding, to address congenital underdevelopment of the breasts, and as part of breast reconstruction after mastectomy for cancer. Implants are medical devices, typically filled with silicone gel or saline and enclosed in a silicone shell, and they are manufactured in a wide range of shapes, sizes, and surface textures. Implants are not expected to last a lifetime; most women will require further surgery at some point to replace, remove, or revise an implant. Careful counselling about realistic expectations, risks, and the need for long-term surveillance is therefore an essential part of the procedure.

\section*{Epidemiology}
Breast augmentation is one of the most commonly performed cosmetic surgical procedures in the world, and the number of procedures has grown steadily over recent decades. It is performed overwhelmingly in women, most commonly in their 20s to 40s, and the prevalence of implants varies widely between countries according to cultural attitudes, cost, and access. Beyond cosmetic augmentation, implants and tissue expanders are widely used for reconstruction after mastectomy, and this application accounts for a substantial proportion of all implant placements. The long-term burden of implants is significant: studies suggest that a large proportion of women who have implants will undergo further surgery within ten years, most commonly for capsular contracture or other implant-related problems.

\section*{Aetiology and Pathophysiology}
Breast augmentation is not a disease, and the reasons for surgery are varied:

\begin{itemize}
    \item \textbf{Cosmetic:} to increase breast size or improve breast shape and projection
    \item \textbf{Asymmetry:} to correct breasts that differ noticeably in size or shape
    \item \textbf{Congenital:} to address underdevelopment of the breasts (hypoplasia), including in conditions such as Poland syndrome
    \item \textbf{Post-pregnancy or weight loss:} to restore lost breast volume and fullness
    \item \textbf{Reconstruction:} to rebuild a breast after mastectomy, either immediately or later
\end{itemize}

After placement, the body responds by forming a thin, fibrous capsule of scar tissue around the implant, which normally holds it securely in place. In some women this capsule progressively tightens and hardens -- capsular contracture. Implants can also rupture, leak, fold, or shift over time, and the scar capsule can become infected or chronically inflamed. A rare cancer of the immune cells lining the capsule, breast implant-associated anaplastic large cell lymphoma (BIA-ALCL), has been linked mainly to textured implants and typically presents with late swelling of the breast.

\section*{Clinical Features}
This section describes what to expect after surgery and what may go wrong, rather than the symptoms of a disease:

\begin{itemize}
    \item \textbf{Normal recovery:} swelling, bruising, and soreness for the first few weeks, settling gradually with time
    \item \textbf{Normal healing:} initial firmness that softens as swelling resolves
    \item \textbf{Capsular contracture:} the breast becomes progressively hard, painful, or visibly distorted, sometimes years after surgery
    \item \textbf{Rupture or leak:} with silicone implants this is often silent and found on imaging; a saline implant deflates noticeably over days
    \item \textbf{Implant movement:} shifting, rotation, or an abnormal appearance of the breast
    \item \textbf{Infection:} redness, warmth, pain, fever, or discharge from the wound in the early weeks
    \item \textbf{Late swelling:} sudden enlargement of one breast years after surgery, which should be investigated promptly because of the link to BIA-ALCL
\end{itemize}

\section*{Differential Diagnosis}
When a problem arises, other causes must be considered:

\begin{itemize}
    \item Capsular contracture versus the normal firmness of early healing
    \item Haematoma or seroma (a fluid collection) versus infection
    \item Implant rupture versus a benign or malignant breast lump
    \item A late fluid collection around an implant versus BIA-ALCL
    \item Breast cancer, which must always be considered, and for which women with implants should continue to be screened
\end{itemize}

\section*{When to Seek Medical Care}
Contact your surgeon or a doctor promptly if you have signs of infection, a sudden change in breast shape or size, a new hard lump, unusual pain, or swelling around an implant. Late swelling of a breast that has been stable for years warrants urgent investigation.

\textbf{Call 000 (or local emergency services) for:}
\begin{itemize}
    \item High fever, chills, and a hot, red, rapidly worsening breast -- possible severe infection
    \item Sudden severe pain or swelling in the breast
    \item Wound breakdown with heavy bleeding or discharge
    \item Difficulty breathing, calf pain, or other symptoms suggesting a blood clot after surgery
\end{itemize}

\section*{Diagnosis and Investigations}
\begin{itemize}
    \item \textbf{Clinical examination:} assessment of the implant, the surrounding scar capsule, and the breast tissue itself
    \item \textbf{Ultrasound:} useful for examining the implant and the surrounding tissue and fluid
    \item \textbf{MRI:} the most reliable imaging test for detecting silicone implant rupture, whether it is symptomatic or silent
    \item \textbf{Mammography:} performed with special displacement views so the implant does not obscure breast tissue, and used for routine screening
    \item \textbf{Aspiration:} fluid from a late collection around an implant can be examined and tested, and its analysis is the way BIA-ALCL is diagnosed
    \item \textbf{Biopsy:} of any suspicious lump, thickened capsule, or abnormal tissue as required
\end{itemize}

\section*{Management}
\begin{itemize}
    \item \textbf{Informed consent:} a thorough discussion with a qualified plastic surgeon of the options, the realistic outcomes, the risks, and the likely need for future surgery
    \item \textbf{Choosing the implant:} decisions about size, shape, surface texture, and position (above or below the pectoral muscle), guided by the woman's anatomy and goals
    \item \textbf{Post-operative care:} wound care, analgesia, and a graduated recovery plan as directed by the surgeon
    \item \textbf{Capsular contracture:} mild cases may be treated conservatively, but established contracture is usually treated surgically by capsulectomy, often with implant replacement or removal
    \item \textbf{Rupture:} the implant is usually removed and, if desired, replaced, with removal of free silicone in the surrounding tissue
    \item \textbf{Infection:} antibiotics, with removal of the implant in persistent or severe cases
    \item \textbf{Reconstruction planning:} in breast cancer, implant-based reconstruction is coordinated with the oncology team, and radiotherapy decisions are central to planning
    \item \textbf{Ongoing surveillance:} regular check-ups of the implants and continued participation in breast screening are recommended for all women with implants
\end{itemize}

\section*{Complications}
\begin{itemize}
    \item Capsular contracture, the most common reason for further surgery
    \item Implant rupture or leaking, the likelihood of which rises with time in situ
    \item Infection, haematoma, or seroma after surgery
    \item Scarring, rippling of the skin, or a result that looks or feels unnatural
    \item Implant displacement, rotation, or folding
    \item Altered or lost sensation in the nipple and breast
    \item The need for further operations over the lifetime of the implant
    \item Breast implant-associated anaplastic large cell lymphoma (BIA-ALCL), a rare but serious cancer linked mainly to textured implants
    \item Possible effects on the accuracy of breast cancer screening and the interpretation of mammograms
\end{itemize}

\section*{Prognosis and Outcomes}
Most women who have breast augmentation are satisfied with the result, and serious complications are uncommon in the short term. However, implants are not permanent devices: on average they last around ten to twenty years, and a substantial proportion of women will eventually need further surgery to replace or remove them, or to treat a complication such as capsular contracture. Regular check-ups and imaging detect problems early, and prompt investigation of any new swelling or lump is important. Women with implants should continue routine breast cancer screening and report any new or unusual breast symptoms to their doctor.

\section*{Prevention and Self-Care}
\begin{itemize}
    \item Choose a qualified, experienced plastic surgeon and discuss all risks and options thoroughly before surgery
    \item Stop smoking and maintain a healthy weight before surgery to optimise healing
    \item Follow post-operative care and wound instructions carefully
    \item Attend regular check-ups of the implants and continue routine breast screening
    \item Know the signs of rupture, infection, and late swelling, and seek review early if they occur
    \item Discuss the risks of textured implants, including BIA-ALCL, before choosing them
    \item Do not delay seeking advice about any new lump, swelling, or change in the breast
\end{itemize}

\section*{Sources and Further Reading}
The following organisations provide reliable, evidence-based information for patients, students, and clinicians.

\begin{itemize}
    \item NHS. \textit{Breast implants and breast augmentation: information on risks and aftercare.} https://www.nhs.uk/conditions/
    \item Mayo Clinic. \textit{Breast augmentation: procedure overview and patient education.} https://www.mayoclinic.org/diseases-conditions/
    \item MedlinePlus. \textit{Breast implants: risks, surveillance, and health information.} https://medlineplus.gov/
    \item MSD Manual Professional Version. \textit{Complications of breast implants and their management.} https://www.msdmanuals.com/
    \item World Health Organization. \textit{Information on implantable medical devices and patient safety.} https://www.who.int/
    \item Centers for Disease Control and Prevention. \textit{Breast implant safety and BIA-ALCL information.} https://www.cdc.gov/
\end{itemize}
